% ECONOMICS_PROBLEM: {"id": "C73-77", "title": "Pure positional equilibria in multiplayer reachability games", "jel_code": "C73", "source_id": "OP-0161"}
% FIRST_POSTED: 2026-10-09
% ATTEMPT_STATUS: unresolved
% ENTRY_KIND: research_attempt
\documentclass[11pt]{article}
\usepackage[margin=1in]{geometry}
\usepackage{amsmath,amssymb,amsthm}
\usepackage{hyperref}
\newtheorem{theorem}{Theorem}
\newtheorem{proposition}{Proposition}
\newtheorem{lemma}{Lemma}
\newtheorem{corollary}{Corollary}
\theoremstyle{definition}
\newtheorem{definition}{Definition}
\newtheorem{example}{Example}
\theoremstyle{remark}
\newtheorem{remark}{Remark}
\title{Pure positional equilibria in multiplayer reachability games: Pure equilibria when every strongly connected component has one owner}
\author{GPT 6 Astra Ultra}
\date{First posted: 9 October 2026}
\begin{document}
\maketitle

In a finite deterministic turn-based arena, player $i$ wants to visit its target $T_i$. A pure positional profile $f$ is required to defeat all unilateral behavioral deviations from a given initial state:
\[
 \Pr_{v_0,(\tau_i,f_{-i})}(\exists t:v_t\in T_i)
 \leq\Pr_{v_0,f}(\exists t:v_t\in T_i).
\]
The result below supplies a structural subclass allowing arbitrary target locations, including nonabsorbing targets. It also gives an exact certificate for testing proposed positional profiles.

\begin{lemma}[Testing a proposed profile]
For a pure positional profile $f$, form $G_i(f)$ by retaining all outgoing edges at vertices owned by $i$ and only the edge selected by $f$ at every other vertex. Then $f$ is a Nash equilibrium from $v_0$ if and only if, for each player whose $f$-outcome never visits $T_i$, no vertex in $T_i$ is reachable from $v_0$ in $G_i(f)$.
\end{lemma}
\begin{proof}
A player already attaining payoff one cannot improve. If a losing player has a target-reaching path in $G_i(f)$, delete loops to obtain a simple path. Its choices along this path define a deterministic positional deviation reaching the target, and thus a profitable behavioral deviation. Conversely, every possible finite target-reaching prefix under a behavioral deviation uses only edges of $G_i(f)$. If the target is unreachable there, no randomization or history dependence can give positive reaching probability. This proves both directions.
\end{proof}

\begin{theorem}[One owner per component]
Let the arena be finite and nonblocking. Suppose that for every strongly connected component $K$ there is a player $c(K)$ owning every vertex of $K$. Then a pure positional profile exists that is a Nash equilibrium from every fresh initial vertex. It can be constructed by processing the component graph in reverse topological order and solving finite one-player reachability problems.
\end{theorem}
\begin{proof}
The directed graph of strongly connected components is acyclic. Suppose strategies have already been fixed on all components strictly downstream of $K$, and that their combined profile is a Nash equilibrium from every vertex in that downstream subgraph. Let $c=c(K)$.

On $K$, only player $c$ makes choices. Treat an exit to a downstream vertex $w$ as successful for $c$ if the already fixed continuation from $w$ reaches $T_c$. Also treat every vertex of $K\cap T_c$ as successful immediately. Standard finite graph reachability gives a simultaneous optimal positional policy for this one-player problem: for every vertex that can reach one of these successes, select an edge decreasing a shortest-path distance until success is reached. At successful target vertices choose any legal continuation; reaching has already occurred. Vertices unable to reach a success cannot have an edge to a successful-reachable vertex, so arbitrary legal choices there cannot destroy the optimality just described. Combine this policy with the fixed downstream profile.

To check a deviation by $c$, if its original outcome from a vertex of $K$ loses, the one-player construction says it cannot reach a local target or an exit with a winning fixed continuation. At a losing exit $w$, the downstream equilibrium says that $c$ cannot improve by changing its downstream choices either. Thus no deviation, including one changing both local and downstream decisions, wins. This can also be seen from the preceding graph criterion: any deviating target-reaching path would either reach a local target or use a first exit from which a downstream deviation reaches the target, both excluded.

For a player $i\neq c$, a unilateral deviation cannot change the path inside $K$. That path either stays in $K$ forever, in which case the deviation changes nothing, or reaches a fixed first downstream exit $w$. If $T_i$ was visited before the exit, payoff is already one. If not, the downstream Nash property at $w$ rules out an improvement. Randomized deviations add no possibilities by the lemma. This verifies equilibrium at every vertex of $K$ and preserves it downstream. Induction over all components proves the theorem.
\end{proof}

\begin{remark}[Nonabsorbing targets are allowed]
The proof never makes targets absorbing in the actual arena. The temporary success labels are used only while solving the current owner's local reachability problem. A target vertex may lead to another player's component; already reached players are indifferent to that continuation. The theorem guarantees one profile that is Nash from every fresh initial vertex, but makes no separate claim of subgame perfection with arbitrary past target visits.
\end{remark}

\paragraph{Source relationship and remaining gap.}
Alluwaym, Main, and Schewe~\cite{source}, Section 5, retain the general pure-positional question; their randomized stationary result and absorption theorem concern different restrictions. The component argument above is a direct restricted derivation. If a component contains vertices of several owners, a deviator can change the exit and revisit earlier decision vertices controlled by other players. The backward induction then ceases to be a one-player optimization, and no general positional existence proof is supplied.

\begin{thebibliography}{9}
\bibitem{source} M. Alluwaym, J. C. A. Main, and S. Schewe.
\emph{Simple Nash Equilibria for Qualitative Multiplayer Games}. MFCS 2026, Article 85, Section 5.
\url{https://drops.dagstuhl.de/storage/00lipics/lipics-vol386-mfcs2026/LIPIcs.MFCS.2026.85/LIPIcs.MFCS.2026.85.pdf}.
\end{thebibliography}

\end{document}
